\section*{Chapitre \ref{ch:mx:auction-mechanics-and-theory} --- Mécanique et théorie des enchères}

\begin{solution}{exo:mx:auction-mechanics-and-theory:1}
À 10.00, la demande est de 900 et l'offre de 500~; à 10.01 et 10.02, la demande est de 500 et l'offre de 800~; à 10.03, la demande est de 300 et l'offre de 1\,000. Le volume, 500, est maximal à
10.00, 10.01 et 10.02~; le surplus est de 400 à l'achat à 10.00 et de 300 à la vente aux autres, donc 10.01 et 10.02 subsistent~; le surplus est du côté vendeur, donc le
plus bas, 10.01~: 500 actions, 300 intentions de vente non exécutées.
\end{solution}

\begin{solution}{exo:mx:auction-mechanics-and-theory:2}
À 50.10, il y a 3\,000 actions de plus à l'achat que de vendeurs~: un vendeur de jusqu'à 3\,000 actions ferait baisser le prix ou serait exécuté à ce prix. Un fournisseur de liquidité qui valorise le titre
en dessous de 50.10 offre une partie des 3\,000 actions à sa valeur majorée d'une prime~; s'il offre à 50.10 ou en dessous, le cours indicatif tend vers son prix.
\end{solution}

\begin{solution}{exo:mx:auction-mechanics-and-theory:3}
Environ $0.15/1=15~\%$ et $0.15/10=1.5~\%$~; avec un écart d'une milliseconde et des lots toutes les 10 millisecondes, environ 10~\%. La simulation a donné 11.75~\%, 1.25~\% et 11.0~\%.
\end{solution}

\begin{solution}{exo:mx:auction-mechanics-and-theory:4}
$3\,000/500=6$ ticks. La simulation a mesuré 4.7 ticks pour un ordre passé à la dernière seconde~: il déplace la clôture par rapport à l'endroit où la séance allait se terminer, et une partie de
l'intérêt affiché qu'il parcourt était déjà apparié à d'autres ordres ou compensé par des contreparties antérieures dans la phase d'enchère.
\end{solution}

\begin{solution}{exo:mx:auction-mechanics-and-theory:5}
La plus grande partie des actions appariées provient de l'intérêt affiché rencontrant le flux précoce d'\omterm{def:mx:the-limit-order-book:orders}{ordres au marché} à la clôture~; les ordres ultérieurs modifient peu l'appariement, mais déterminent quel
ordre marginal fixe le prix, si bien que le cours indicatif continue de bouger alors que le volume apparié varie à peine.
\end{solution}

\begin{solution}{exo:mx:auction-mechanics-and-theory:6}
Personne ne sait si un ordre envoyé à l'heure prévue sera le dernier~: des contreparties peuvent encore arriver, de sorte que l'effet d'un ordre tardif sur le prix peut
être contrebalancé (4.7 ticks deviennent 3.7 dans la simulation). Les participants qui doivent être présents à l'enchère doivent entrer avant la fenêtre et supporter l'incertitude
sur l'heure exacte de leur exécution.
\end{solution}

\begin{solution}{exo:mx:auction-mechanics-and-theory:7}
1~ms~: 11.75~\% contre 15~\%~; 10~ms~: 1.25~\% contre 1.5~\%~; 100~ms~: 0.25~\% contre 0.15~\%. L'approximation suppose que les sauts se répartissent uniformément à l'intérieur des lots
et qu'une seule frontière de lot décide~; avec 400 sauts, les taux sont estimés à environ un point de pourcentage près, et lorsque l'écart dépasse l'intervalle, le taux
vaut un.
\end{solution}

\begin{solution}{exo:mx:auction-mechanics-and-theory:8}
La clôture est fixée par les ordres marginaux, et le dernier entré n'a personne pour y répondre~: dans la simulation, un achat de 3\,000 actions passé à la dernière seconde d'une
enchère à fin fixe a déplacé la clôture de 4.7 ticks, soit 1.75 de plus que le même ordre passé une minute plus tôt. Les heures limites, les collars et les \omterm{def:mx:auction-mechanics-and-theory:randomend}{fins aléatoires} existent parce qu'un seul
ordre tardif peut la déplacer.
\end{solution}

\begin{solution}{pb:mx:auction-mechanics-and-theory:1}
\textbf{1.} Volume maximal, surplus minimal, pression du marché, proximité de la référence~; 10.01, 500 actions.
\textbf{2.} Un cours indicatif, des actions appariées et un déséquilibre, à chaque seconde de l'enchère simulée, dès qu'une partie du volume est appariée.
\textbf{3.} Un intérêt affiché de 500 actions par tick, des \omterm{def:mx:the-limit-order-book:orders}{ordres au marché} à la clôture aléatoires, un achat de 10\,000 actions d'un fonds et des contreparties.
\textbf{4.} À 100.01 ou 99.99, leur propre valeur plus un tick, en proportion de la distance dont l'indicatif a été repoussé~; offrir au prix repoussé
figerait le mouvement.
\textbf{5.} 4.7, 3.5, 2.2 et 1.0 ticks.
\textbf{6.} Plus aucun ordre n'arrive dans les cinq dernières secondes.
\textbf{7.} 76~\% et 89~\%.
\textbf{8.} Nasdaq~: information sur le déséquilibre entre 15~h~50 et 16~h~; le LSE~: une période aléatoire pouvant atteindre 30 secondes après chaque enchère et chaque prolongation.
\textbf{9.} Voir les définitions.
\textbf{10.} 3.0, 3.5 et 4.7 ticks.
\textbf{11.} 3.0, 3.5 et 3.7 ticks.
\textbf{12.} 1.75 tick avec une fin fixe, 0.75 avec une \omterm{def:mx:auction-mechanics-and-theory:randomend}{fin aléatoire}.
\textbf{13.} Celui qui veut déplacer la clôture y gagne~; l'autre côté de l'enchère paie le prix déplacé.
\textbf{14.} La mise en relation continue fait de la vitesse le facteur qui décide qui prend les cotations périmées~; des enchères par lots toutes les fractions de seconde suppriment la valeur des petits
avantages de vitesse.
\textbf{15.} 99.75~\%, 11.75~\%, 1.25~\% et 0.25~\% pour un écart de 150 microsecondes à 0, 1, 10 et 100 ms~; 99.75~\%, 99.75~\%, 11.0~\% et 0.75~\% pour une milliseconde.
\textbf{16.} Le sniper ne gagne que si un lot se clôt entre son ordre et l'annulation du teneur de cotation~: avec un instant de saut uniforme, la probabilité vaut écart$/$intervalle.
\textbf{17.} Attendre jusqu'à un intervalle pour chaque exécution, et disposer de moins d'informations à partir du cours continu.
\textbf{18.} \emph{Résultat nommé}~: l'indicatif se situe à 4.7 ticks du cours final deux minutes avant le croisement, à 1.0 tick à dix secondes et il est exact dans les dernières
cinq secondes~; passer un achat de 3\,000 actions à la dernière seconde plutôt qu'une minute plus tôt déplace la clôture de 1.75 tick de plus avec une fin fixe et de 0.75 avec une \omterm{def:mx:auction-mechanics-and-theory:randomend}{fin
aléatoire} pouvant atteindre 30 secondes.
\textbf{19.} Oui~: cela réduit de plus de moitié la valeur d'arriver en dernier, au prix d'une heure de fin incertaine.
\textbf{20.} Des lots plus longs que les écarts de latence qui décident des courses, environ une milliseconde ou plus pour les écarts simulés ici, suppriment l'essentiel du sniping~;
l'intervalle doit rester assez court pour les traders qui veulent une exécution immédiate.
\end{solution}

\begin{solution}{iq:mx:auction-mechanics-and-theory:1}
Tous les ordres sont collectés et compensés à un seul prix qui maximise le volume exécuté (puis minimise le surplus, suit son côté, puis la référence)~;
le déséquilibre est la quantité restant d'un côté au cours indicatif.
\iqlookfor{Prix uniforme~; l'ordre des règles~; le déséquilibre.}
\end{solution}

\begin{solution}{iq:mx:auction-mechanics-and-theory:2}
Suffisamment tôt pour être avant l'heure limite et laisser l'intérêt de contrepartie réagir au déséquilibre publié~; je fractionne l'ordre si sa taille est grande par rapport au volume
de clôture habituel~; je surveille les publications de déséquilibre, l'écart entre l'indicatif et le cours continu et l'arrivée de contreparties.
\iqlookfor{Heures limites~; laisser l'autre côté réagir~; suivi du déséquilibre.}
\end{solution}

\begin{solution}{iq:mx:auction-mechanics-and-theory:3}
Pour supprimer la certitude d'arriver en dernier et donc l'incitation à attendre pour déplacer le prix. Mesurer la part du volume et des variations de prix dans les dernières
secondes, ainsi que l'écart entre le dernier indicatif et le prix final, avant et après la réforme, ou par rapport à un marché comparable sans elle.
\iqlookfor{Incitation à être le dernier~; activité tardive et écarts comme mesures.}
\end{solution}

\begin{solution}{iq:mx:auction-mechanics-and-theory:4}
Des adjudications à prix uniforme à intervalles courts remplacent la mise en relation continue~: les ordres d'un même intervalle sont simultanés, de sorte que le trader le plus rapide
ne prend plus toutes les cotations périmées~; le sniping tombe à environ l'écart de latence divisé par l'intervalle.
\iqlookfor{La course~; temps discret~; l'intuition écart sur intervalle.}
\end{solution}

\begin{solution}{iq:mx:auction-mechanics-and-theory:5}
Demande et offre cumulées par \omterm{def:mx:the-limit-order-book:tick}{niveau de prix} (deux tableaux triés ou arbres de tailles agrégées, les \omterm{def:mx:the-limit-order-book:orders}{ordres au marché} étant traités à part), mises à jour en $O(\log n)$ par
ordre~; l'indicatif recherche le maximum de $\min(D,S)$ avec les règles de départage, ou le conserve de façon incrémentale près du dernier indicatif.
\iqlookfor{Profondeur agrégée~; sommes cumulées~; départage.}
\end{solution}

\begin{solution}{iq:mx:auction-mechanics-and-theory:6}
Comparer ses ordres avec ses positions et avec les instruments valorisés à partir de la clôture (produits dérivés, valorisations de fonds), mesurer l'impact sur le prix de ses
ordres tardifs et vérifier si les cours reviennent après la clôture~; les ordres tardifs légitimes ont une raison de négocier à la clôture et ne présentent pas un tel schéma.
\iqlookfor{Preuves d'intention~; retour des cours après la clôture~; positions liées.}
\end{solution}
